\noindent\textbf{(a)} $f(x)=x/2+C/x$ en un intervalo donde $x\ne0$.

\noindent\textbf{(b)} $x^4+2x^2y+4y=K$, es decir $y=(K-x^4)/(2x^2+4)$, en intervalos con $x\ne0$.

\noindent\textbf{(c)} $e^{-4y}(x^2+y+1/4)=K$, en sus arcos regulares. La ortogonalidad respecto de las soluciones gráficas originales se afirma en $x\ne0$; la forma diferencial se extiende a $x=0$ salvo la dirección indeterminada en $(0,0)$.
